\section{Discussion}\label{sec:discussion}

\paragraph{What makes ICS explanations timely?}
Our results decompose the timeliness problem into three factors.
\begin{enumerate}
\item \emph{Detection dominates.} On PCL-Sim, \intimeFailUndeliveredPct\,\% of the deadline failures of the best configurations are attacks detected too late or not at all. No explanation method can fix those, so the deadline perspective puts detection latency \citep{puccetti2026latency} back at the centre.
\item \emph{When attributions are computed matters more than how.} Every attribution method we tested, gradient-based, sampling-based or exact, degrades as an attack propagates through coupled loops. Explanations computed a minute or more after the alarm are therefore often misleading if read in isolation. This is consistent with, and adds a mechanism to, the timing failure reported by \citet{fung2024attributions}.
\item \emph{Temporal evidence accumulation is a cheap and general remedy.} It adds no detector evaluations, and it applies unchanged to any attribution method. In our experiments it was the single largest improvement for every method. It also turns flickering top-1 outputs into stable ones.
\end{enumerate}

\paragraph{What does control-loop structure buy?}
The structure controls answer a question left open by grouped Shapley methods \citep{owen1977unions,franco2026shats}: grouping is only useful if the groups match how the plant is controlled. Correlation clusters and random partitions with the same group sizes perform dramatically worse than control loops.

Given loop structure, exact Shapley values (\LX) are the most accurate black-box attribution at the alarm on HAI, and tie with the best methods once \TEA{} is applied. Their operational advantage is predictability, not a large accuracy margin:
\begin{itemize}
\item a fixed number of detector queries ($2^L$), decided at design time and independent of the data;
\item no randomness, so the same alarm yields the same explanation;
\item no gradient access.
\end{itemize}
Where gradients are available, integrated gradients aggregated per loop with \TEA{} is an equally strong and cheaper choice in our data. We recommend reporting it as a baseline in future work. The training-background variant (\LX-bg) was slightly better on the simulator and slightly worse at the alarm on HAI. We therefore do not claim a general advantage for the incident-local reference.

\paragraph{Practical recommendations}
\begin{enumerate}
\item Evaluate explanations at operational timing, against documented targets, and over time, not only at a single window.
\item Group attributions by control loop, using the P\&ID or controller configuration.
\item Accumulate evidence across the incident and present a short ranked list (the top three contained the attacked loop in 94--95\,\% of detected attacks at 60\,s in both benchmarks).
\item Report deadline-conditioned metrics alongside latency, because fast explanations of late or wrong alarms do not help the operator.
\item Use chronological, attack-disjoint evaluation. On HAI, a random window split inflated F1 from \leakBlockedFOne{} to \leakRandomFOne.
\end{enumerate}

\subsection{Limitations and threats to validity}\label{sec:limitations}
\paragraph{External validity}
\begin{itemize}
\item \emph{Benchmark coverage.} We evaluate on one real HIL testbed (HAI 21.03, 50 attacks) and one simulator of our own design. SWaT and WADI, the most common ICS benchmarks, are only available on request, and we could not include them in this study. Replication on them, on HAI 22.04/23.05, and on TEP-based manipulation suites \citep{fung2024attributions} is the most important next step.
\item \emph{Other CPS domains.} Both benchmarks are industrial process CPS. Power grids, vehicles and medical devices also have control loops and physical deadlines, so the framework applies to them in principle. Their loop sizes, sampling rates and deadlines differ, however, and we make no empirical claim about them.
\item \emph{Simulator bias.} A simulator designed by the authors of a method risks favouring it. We mitigate this in three ways: all explainers query the same score function and reference; attack primitives mirror HAI's taxonomy; and simulator parameters were fixed on development seeds before any explainer was evaluated. The simulator is released for scrutiny.
\end{itemize}

\paragraph{Deadlines} Hazard deadlines are computed in simulation, with no operator or safety-system response, and with safety limits we chose. On HAI we can only report DCL over a deadline grid, because recorded data carry no process safety times. Real process safety times come from hazard analyses (IEC~61511) and differ by orders of magnitude between processes.

\paragraph{Hardware} Latencies were measured on one thread of a server CPU, not on a PLC-adjacent embedded device. We report hardware-independent evaluation counts and a sensitivity analysis with an assumed slowdown factor, but we did not measure on embedded hardware. Edge-class conclusions are therefore extrapolations.

\paragraph{Method assumptions}
\begin{itemize}
\item \emph{Engineering metadata.} \LX{} needs a loop map. Some HAI assignments are our engineering judgement (flagged in the released configuration).
\item \emph{Scaling.} Exact computation is exponential in $L$. Beyond $L\approx 14$ (16k evaluations), loops would need a two-level hierarchy (unit $\to$ loop) or sampling, and the deterministic-latency guarantee would weaken.
\item \emph{Reference.} The incident-local reference assumes that the plant was in a normal state $\Delta_{\mathrm{ref}}$ seconds before the alarm. Slow ramps that begin long before detection contaminate it (\cref{sec:ablation}).
\item \emph{Detector family.} We study forecasting detectors only. Reconstruction-, density- or classifier-based detectors would need their own score functions.
\end{itemize}

\paragraph{Adversary} The decoy experiment considers a non-adaptive adversary with a fixed small bias. Explanation-manipulation attacks that preserve the detector's decision while steering SHAP, LIME or integrated-gradients explanations have been demonstrated on security classifiers, including intrusion detection \citep{slack2020fooling,mia2025vulnerable}. Whether such an adversary can steer a loop-level attribution while respecting the physical constraints of a closed control loop remains open and is the most important robustness question for this work.

\paragraph{Human factors} We measure whether the correct loop is ranked first or within the top $k$. We do not measure whether operators act faster or better. The preference for showing several candidates found in the operator survey of \citet{fung2024attributions} motivates our HR@3 results but does not substitute for a user study.

\paragraph{Literature search} Bibliographic APIs (OpenAlex, Crossref, arXiv, Semantic Scholar) were not reachable from our analysis environment. The novelty check therefore combined the curated corpus with general web search. We found no prior work combining the elements above, but this does not prove that none exists.

